\documentclass[a4paper,11pt]{article}
\usepackage[margin=23mm]{geometry}
\IfFileExists{../base/natalreport.sty}{\usepackage{../base/natalreport}}{\usepackage{../natalreport}}
\usepackage{booktabs,longtable}
\directlua{dofile("dados-relatorio.lua")}
\NatalConfig{chart radius=0.38\linewidth}
\title{Mapa natal\\\large Relatório técnico}
\author{}
\date{}
\begin{document}
\maketitle
\section{Roda do mapa}
\begin{center}\NatalChartInternal\end{center}
\clearpage
\section{Posições}
R indica retrogradação na entrada.
\begin{longtable}{@{}llrcc@{}}
\toprule Corpo ou ponto & Signo & Posição & Casa & R\\ \midrule
\endhead
\NatalPosicoes
\bottomrule
\end{longtable}

\section{Cúspides}
\begin{longtable}{@{}clrl@{}}
\toprule Casa & Signo & Posição & Ângulo\\ \midrule
\endhead
\NatalCusps
\bottomrule
\end{longtable}

\clearpage
\section{Grade de aspectos}
A numeração das colunas segue a ordem das linhas. Cj é conjunção, Sx é sextil,
Qd é quadratura, Tr é trígono e Op é oposição.
\begin{center}
\small\setlength{\tabcolsep}{3pt}
\begin{tabular}{@{}l*{15}{c}@{}}
\toprule
Corpo ou ponto & 1 & 2 & 3 & 4 & 5 & 6 & 7 & 8 & 9 & 10 & 11 & 12 & 13 & 14 & 15\\
\midrule
\NatalGrade
\bottomrule
\end{tabular}
\end{center}

\section{Lista completa de aspectos}
A lista inclui as relações com os ângulos.
\begin{longtable}{@{}lllrl@{}}
\toprule Origem & Destino & Tipo & Orbe & Fase\\ \midrule
\endhead
\NatalAspectos
\bottomrule
\end{longtable}

\section{Velocidades informadas}
A direção vem do sinal do delta. O rótulo e a razão reproduzem os valores do arquivo.
\begin{longtable}{@{}lrlrr@{}}
\toprule Corpo & Delta & Direção & Rótulo & Razão\\ \midrule
\endhead
\NatalVelocidades
\bottomrule
\end{longtable}
\end{document}